\subsection{Three further vanishing degrees and the isolated residual prime}
\label{sec:higher-k-vanishing}

\begin{theorem}[K17-K16-K20-K24]\label{thm:higher-k-vanish}
The stipulated arithmetic theorem of Alp\"oge implies
\[
 K_{16}(\mathbf Z)=K_{20}(\mathbf Z)=K_{24}(\mathbf Z)=0.
\]
Together with the previous degree-twelve theorem and the known degrees
four and eight, this gives
\[
                  K_{4k}(\mathbf Z)=0\qquad(1\le k\le6).
 \tag{HK.1}
\]
For every odd prime $\ell$, the cyclotomic class components
$A_\ell^{[-8]}$, $A_\ell^{[-10]}$, and $A_\ell^{[-12]}$ vanish.
The full-domain Soul\'e characters $\kappa_9$, $\kappa_{11}$, and
$\kappa_{13}$ are surjective onto $\mathbf Z_\ell$.
\end{theorem}

\begin{proof}
Take respectively $e=8,10,12$, so the desired group is $K_{2e}$.
The source datum uses the same auxiliary field
$E=\mathbf Q(\sqrt{-3})$, modulus two and trivial character as before.
For every prime $\ell>2^{31}$ all of its hypotheses hold: $e$ is
positive and even, $\ell\ne2,3$, and $\ell-1>2e$.
Corollary~\ref{cor:u3-general-coverage} supplies every rank-three
source system in the complete coefficient space of
Theorem~\ref{thm:u3-general}. Lemma~\ref{lem:gl2-conductor} supplies
every rank-two system in the complete classical space of
Proposition~\ref{prop:gl2-general}. All weight and Frobenius
normalizations are exactly those used by the general unit criterion.

For $e=8$ and $12$, (G7) shows that
\[
               \gcd(N_{e,7},N_{e,13})
\]
has no prime factor at least $2^{31}$. Thus, for every prime $\ell$
in the present range, at least one of $N_{e,7}$ and $N_{e,13}$ is
an $\ell$-adic unit. Choose that auxiliary prime $p$; it is automatically
different from $\ell$ and split in $E$.
Corollary~\ref{cor:general-hecke-unit} makes the strict negative Selmer
length zero. By (K5), $K_{2e}(\mathbf Z)\{\ell\}=0$.

For $e=10$, the same argument works at every prime $\ell>2^{31}$
except possibly
\[
                         q=9496386083.
\]
This is the only large prime factor in (G7). At that prime the separate
exact computation in Proposition~\ref{prop:bernoulli-residue} gives
\[
 B_{q-11}\equiv3753150554\not\equiv0\pmod q.
\]
The index satisfies $2\le q-11\le q-3$, and its denominator is prime
to $q$. Lemma~\ref{lem:isolated-reflection}, with $e=10$, yields
$A_q^{[11]}=0$, then $A_q^{[-10]}=0$, and finally
$K_{20}(\mathbf Z)\{q\}=0$ by (K2).
This verifies one specified component at $q$; it does not assert
Vandiver for all components of that prime.

For every odd prime $\ell<2^{31}$ the already cited published
Vandiver verification gives all three primary vanishings, and
Rognes--Weibel gives the 2-primary vanishings. The groups $K_{2e}$
are finite, so they are zero. Combining with
Theorem~\ref{thm:k12-vanish} and (K4) proves (HK.1).
The class-component assertions are (K2) at $k=4,5,6$.
The Soul\'e image-index formula (K18), applied at $m=e+1=9,11,13$,
gives index one and proves the asserted full-domain surjectivity.
\end{proof}

\begin{corollary}[K17-SIX-DEGREES-PERIOD]
\label{cor:six-degrees-period}
For every odd prime $\ell$ and every positive integer $k$ whose residue
modulo $(\ell-1)/2$ is represented by one of
$0,1,2,3,4,5,6$, one has
\[
                         K_{4k}(\mathbf Z)\{\ell\}=0.
\]
For every odd $\ell$, the strict groups
$\mathcal S_\ell(e)$ vanish for $e=2,4,6,8,10,12$.
\end{corollary}
\begin{proof}
For the nonzero listed residues use (K2) and (HK.1), since $-2k$
depends only on the residue of $k$ modulo $(\ell-1)/2$.
The residue zero is the trivial class-group component, already shown
to vanish after (K3). The strict Selmer assertion is exactly (K5).
\end{proof}

\paragraph{Consequences and limits in the programme.}
These are new specified vanishings from the stipulated arithmetic
source, with a separate unconditional component calculation at the
single residual prime in degree twenty. They update the strict
negative-twist rows for $\operatorname{Spec}\mathbf Z[1/\ell]$ in
SA06's E.4 language. They give full-domain odd Soul\'e surjectivity
through weight thirteen, at odd primes. They do not identify the
Kummer tangent torsor of cote~149 with its geometric subset, nor
remove the restricted-character factorial factors or the irregular
integral Lie defects. The all-irregular-prime kernel question still
has its explicitly retained Iwasawa gate.
